\documentclass[11pt]{article}

\usepackage[letterpaper,margin=1in]{geometry}
\usepackage{microtype}
\usepackage{amsmath}
\usepackage{amssymb}
\usepackage{booktabs}
\usepackage{array}
\usepackage{enumitem}
\usepackage{float}
\usepackage{xcolor}
\usepackage{listings}
\usepackage[hidelinks]{hyperref}

\newcommand{\code}[1]{\texttt{#1}}
\newcommand{\rfckw}[1]{\textsc{#1}}

\definecolor{lstbg}{HTML}{F6F8FA}
\definecolor{lstrule}{HTML}{D0D7DE}
\lstdefinestyle{json}{
  basicstyle=\ttfamily\scriptsize,
  backgroundcolor=\color{lstbg},
  frame=single, rulecolor=\color{lstrule},
  framesep=6pt, xleftmargin=6pt, xrightmargin=6pt,
  breaklines=true, breakatwhitespace=true,
  columns=fullflexible, keepspaces=true,
  showstringspaces=false, aboveskip=10pt, belowskip=6pt,
}

\title{\bfseries A CycloneDX VEX Profile for\\ FedRAMP Rev5 VDR/VER Disposition and Response}
\author{Matthew Venne \\ \small Chief Technology Officer, stackArmor}
\date{June 2026}

\begin{document}
\maketitle

\begin{abstract}
\noindent
This is a companion to \emph{A Deterministic, CVSS--Environmental Method for FedRAMP
Rev5 VDR/VER Vulnerability Prioritization}, which assigns every detected finding a
Potential Agency Impact rating (PAIN, N1--N5) and a remediation deadline by formula.
That scoring is deterministic and computable \emph{a priori}. Deciding whether a
finding is \emph{truly} exploitable in the running environment --- and, when a fix will
miss the deadline, attesting the compensating response --- is a separate,
\emph{a posteriori} step that can only follow detection. FedRAMP's Vulnerability
Evaluation Requirements (VER) make that step normative but name no format. This
proposal recommends \textbf{CycloneDX VEX} as the structured, machine-readable, signable
carrier for it, chosen for the fit of its disposition and response \emph{vocabulary} ---
not because it carries scores; it does not. We give the field mapping to VER, a thin
worked example, an OCI-native distribution pattern, a deployment-context service/workload
attestation pattern, and an optional crosswalk to CISA's status-justification labels for
downstream interoperability. This document is an
\emph{implementation profile}: it shows one defensible way to satisfy existing VER
requirements; it does not propose changing them. The key words \rfckw{must},
\rfckw{should}, and \rfckw{may} are to be read as in RFC~2119.
\end{abstract}

\section{Where disposition sits in the lifecycle}
The method memo computes PAIN and a deadline for \emph{every} finding a scanner emits.
But a scan hit is a candidate, not a confirmed exposure: the flagged code may be
absent, may never execute, may require a configuration nobody enabled, or may sit
behind a control that already neutralizes it. Resolving which is the \emph{disposition}
step.

Disposition is an \emph{a posteriori} determination. The factors that decide it ---
reachability, configuration, surrounding mitigations --- are specific to a given
vulnerability on a given deployment and cannot be enumerated before that vulnerability
is known. The analysis therefore necessarily follows detection, acutely so for novel
vulnerabilities, where no prior due-diligence could exist. This is a statement about
\emph{timing}, not about who or what performs it: the determination \rfckw{may} be fully
automated or AI/agentic, but it cannot be precomputed.

Because VER's clocks are short, the determination must be captured in a form that
machines can produce, exchange, sign, and act on. A disposition resolves a finding onto
one of two tracks:
\begin{itemize}[leftmargin=1.5em,nosep,topsep=3pt]
\item \textbf{Not truly affected} --- the finding is suppressed, with a machine-readable
  justification recording \emph{why} (code absent, unreachable, configuration-gated, or
  already mitigated).
\item \textbf{Affected but remediation will be late} --- the finding is accepted within
  process, with an attested mitigation and a planned response, rather than silently
  breaching the deadline.
\end{itemize}

\section{Why CycloneDX: the vocabulary fits the lifecycle}
FedRAMP is not bound to any single VEX format, and VER specifies \emph{fields and
content}, not a schema (Section~\ref{sec:vermap}). The provider therefore chooses the
carrier --- the same latitude FedRAMP itself exercised in substituting PAIN for SSVC.
On that basis the deciding factor is \emph{expressiveness}: which format's disposition
vocabulary most precisely captures the real outcomes above. CycloneDX's \code{analysis}
object is the strongest fit. The same granularity pays off when disposition is
automated: because each justification names a distinct precondition, an automated or
AI/agentic analyzer can map each evidence source one-to-one to a justification, where a
coarser vocabulary would collapse several distinct findings onto a single label.

\begin{table}[H]
\centering
\footnotesize
\renewcommand{\arraystretch}{1.2}
\begin{tabular}{>{\raggedright\arraybackslash}p{0.30\textwidth} >{\raggedright\arraybackslash}p{0.62\textwidth}}
\toprule
\textbf{CycloneDX field} & \textbf{Values (and what they capture)} \\
\midrule
\code{analysis.state} & \code{resolved}, \code{resolved\_with\_pedigree},
  \code{exploitable}, \code{in\_triage}, \code{false\_positive},
  \code{not\_affected} --- note a first-class \code{false\_positive} and an
  \code{in\_triage} (work-in-progress) state. \\
\code{analysis.justification} & \code{code\_not\_present}, \code{code\_not\_reachable},
  \code{requires\_configuration}, \code{requires\_dependency},
  \code{requires\_environment}, \code{protected\_by\_compiler},
  \code{protected\_at\_runtime}, \code{protected\_at\_perimeter},
  \code{protected\_by\_mitigating\_control} --- nine reasons, naming configuration,
  reachability, and perimeter/runtime controls explicitly. \\
\code{analysis.response} & \code{can\_not\_fix}, \code{will\_not\_fix}, \code{update},
  \code{rollback}, \code{workaround\_available} --- the planned action, which the
  late-remediation track needs and lighter formats lack. \\
\code{analysis.detail} & free text for human-readable context and evidence. \\
\bottomrule
\end{tabular}
\caption{CycloneDX's \code{analysis} block: states, justifications, and responses.}
\end{table}

The operational cases a CSP actually encounters map onto native values without
contortion:
\begin{table}[H]
\centering
\footnotesize
\renewcommand{\arraystretch}{1.25}
\begin{tabular}{>{\raggedright\arraybackslash}p{0.42\textwidth} >{\raggedright\arraybackslash}p{0.50\textwidth}}
\toprule
\textbf{Real-world disposition} & \textbf{CycloneDX expression} \\
\midrule
Scanner false positive & \code{state: false\_positive} \\
Vulnerable code absent / removed & \code{not\_affected} + \code{code\_not\_present} \\
Vulnerable path never executes & \code{not\_affected} + \code{code\_not\_reachable} \\
Needs an unset configuration & \code{not\_affected} + \code{requires\_configuration} \\
Neutralized at the edge (WAF/IAP) & \code{not\_affected} + \code{protected\_at\_perimeter} \\
Neutralized at runtime / by a control & \code{not\_affected} + \code{protected\_at\_runtime} / \code{protected\_by\_mitigating\_control} \\
Affected; fix will miss deadline & \code{exploitable} + \code{response:[workaround\_available, update]} + \code{detail} \\
Under evaluation & \code{in\_triage} \\
\bottomrule
\end{tabular}
\caption{Operational dispositions and their native CycloneDX encodings.}
\label{tab:cases}
\end{table}

\section{Mapping to FedRAMP VER}
\label{sec:vermap}
VER requires the disposition step and a machine-readable report, in each case by
function rather than by named format. CycloneDX satisfies each clause:

\begin{table}[H]
\centering
\footnotesize
\renewcommand{\arraystretch}{1.25}
\begin{tabular}{>{\raggedright\arraybackslash}p{0.30\textwidth} >{\raggedright\arraybackslash}p{0.62\textwidth}}
\toprule
\textbf{VER requirement} & \textbf{CycloneDX mechanism} \\
\midrule
\code{VER-EVA-EFP} (evaluate false positives) & \code{analysis.state = false\_positive} \\
\code{VER-EVA-AIA} (assume exploitable absent evidence) & a VEX statement \emph{is} the
  ``evidence otherwise''; its absence leaves the assume-exploitable default in force \\
\code{VER-RPT-VDT} (per-finding disposition) & \code{analysis.state} +
  \code{justification} + \code{response} keyed to the CVE and component \\
\code{VER-TFR-MAV} / \code{VER-RPT-AVI} (accepted vulns) & \code{state = exploitable}
  with \code{response} and \code{detail} recording the acceptance and mitigation \\
\code{VER-TFR-MRH} (machine-readable JSON history) & CycloneDX is JSON; statements
  accrete over time with timestamps \\
\bottomrule
\end{tabular}
\caption{FedRAMP VER clauses and the CycloneDX fields that satisfy them.}
\end{table}

\section{A thin attestation: what it carries, and what it does not}
The disposition artifact stays deliberately small. It carries only what identifies the
finding and records the judgment:
\begin{itemize}[leftmargin=1.5em,nosep,topsep=3pt]
\item the vulnerability \code{id} and its \code{source};
\item the affected subject, by \code{affects[].ref}: either a software component
  \code{bom-ref} (for artifact applicability) or a service/workload \code{bom-ref}
  (for deployment-context reachability);
\item the \code{analysis} block --- \code{state}, \code{justification},
  \code{response}, \code{detail};
\item provenance --- author, timestamp, and a signature over the artifact.
\end{itemize}

It deliberately \emph{omits} PAIN, the LEV/IRV axes, the Certification Class, and the
CVSS vector. Those are computed deterministically by the method and live in the
provider's scoring report; duplicating them into the disposition artifact would invite
drift and tempt a reader to treat an attestation as a place to re-score. The two
records are joined on a stable key --- the \textbf{CVE identifier plus the affected
subject reference} --- so a reporting layer can attach each disposition to the
finding it governs and recover the retained PAIN for audit. The disposition changes
a finding's \emph{handling}; it never changes its score.


\section{Service and workload subjects: reusable surface attestations}
CycloneDX can describe both software \code{components} and network-facing \code{services}.
Its vulnerability \code{affects[].ref} field references a subject by \code{bom-ref}, and
that subject can be a service as well as a package or image component. This matters for
FedRAMP reachability because there are two different facts in play:
\begin{itemize}[leftmargin=1.5em,nosep,topsep=3pt]
\item a \textbf{surface attestation} says what the internet can reach on a deployed
  service or workload --- the exposed-surface set $E(a)$;
\item a \textbf{finding disposition} says whether a particular vulnerability's required
  surface $X(v)$ intersects that exposed surface.
\end{itemize}

A service-level surface attestation is therefore reusable across findings. It does not,
by itself, declare every CVE on the service reachable or unreachable; instead it caches
the asset-side evidence so each CVE can be evaluated automatically without rediscovering
Ingress, Gateway, load-balancer, backend protocol, ALPN, and edge-auth state. If a service
attestation says the subject is not internet-reachable at all, every finding on that subject
can be classified NIRV. If it says the subject is internet-reachable, each \code{AV:N}
finding remains IRV by default unless the per-CVE required surface is known and positively
mismatches the attested exposed surface.

A stable service/workload \code{bom-ref} \rfckw{should} identify the deployment subject,
not merely the image artifact. A Kubernetes service, for example, can be represented with
a stable URN-like reference:
\begin{lstlisting}[style=json]
{
  "services": [
    {
      "bom-ref": "urn:k8s:prod:default:service:payments-api",
      "name": "payments-api",
      "endpoints": ["https://api.example.com"],
      "properties": [
        { "name": "k8s:cluster", "value": "prod" },
        { "name": "k8s:namespace", "value": "default" },
        { "name": "k8s:kind", "value": "Service" },
        { "name": "k8s:name", "value": "payments-api" },
        { "name": "vdr:assetInternetReachable", "value": "true" },
        { "name": "vdr:exposedBackendProtocol", "value": "http" },
        { "name": "vdr:exposedBackendProtocolVersion", "value": "HTTP2" },
        { "name": "vdr:exposedBackendAlpn", "value": "h2" },
        { "name": "vdr:routeEvidence", "value": "HTTPRoute public/payments -> Service default/payments-api:8080" }
      ]
    }
  ]
}
\end{lstlisting}

The per-finding record can then reference the same service subject and record the
automated disposition for a specific CVE:
\begin{lstlisting}[style=json]
{
  "id": "CVE-2026-30303",
  "analysis": {
    "state": "exploitable",
    "detail": "Asset is internet-reachable and CVSS is AV:N. Required protocol is unknown, so the required surface is treated as a wildcard."
  },
  "affects": [{ "ref": "urn:k8s:prod:default:service:payments-api" }],
  "properties": [
    { "name": "vdr:findingInternetReachable", "value": "true" },
    { "name": "vdr:reachabilityDecision", "value": "insufficient_evidence_to_downgrade" },
    { "name": "vdr:remediationTrack", "value": "LEV+IRV" }
  ]
}
\end{lstlisting}

This split preserves automation without overclaiming. The service attestation is reused
for all CVEs on that deployed subject, while the per-CVE disposition still accounts for
\code{AV}, required protocol, ALPN/version, and any vulnerability-specific preconditions.
The surface attestation \rfckw{should} be re-evaluated when the facts that justify it
change: Service selector, selected port, Ingress or Gateway binding, backend protocol,
ALPN policy, edge authentication, or the workload template/image digest when that digest
is part of the evidence. Routine replica churn need not invalidate the attestation when
those evidence inputs remain unchanged.

\section{Worked example}
A single signed CycloneDX document can carry several statements. Below, one finding is
dispositioned \code{not\_affected} (unreachable code) and one is \code{exploitable} with
an attested perimeter mitigation because the vendor fix will arrive after the deadline.
Note the absence of any score, rating, or PAIN field.

\begin{lstlisting}[style=json]
{
  "bomFormat": "CycloneDX",
  "specVersion": "1.6",
  "version": 1,
  "metadata": {
    "timestamp": "2026-06-28T14:00:00Z",
    "authors": [{ "name": "CSP Security Team" }],
    "component": {
      "type": "application",
      "name": "payments-api",
      "bom-ref": "payments-api@2.4.1"
    }
  },
  "vulnerabilities": [
    {
      "id": "CVE-2026-10101",
      "source": { "name": "NVD",
        "url": "https://nvd.nist.gov/vuln/detail/CVE-2026-10101" },
      "analysis": {
        "state": "not_affected",
        "justification": "code_not_reachable",
        "detail": "Vulnerable parser entrypoint is never invoked; only the
                   signing helper from this library is linked."
      },
      "affects": [{ "ref": "pkg:golang/example.com/parser@1.2.3" }]
    },
    {
      "id": "CVE-2026-20202",
      "source": { "name": "NVD",
        "url": "https://nvd.nist.gov/vuln/detail/CVE-2026-20202" },
      "analysis": {
        "state": "exploitable",
        "response": ["workaround_available", "update"],
        "detail": "Vendor fix ETA exceeds the VER deadline. Perimeter WAF
                   virtual patch deployed and egress restricted; tracked for
                   update on vendor release."
      },
      "affects": [{ "ref": "pkg:deb/debian/libxyz@2.1.0-1" }]
    }
  ]
}
\end{lstlisting}

\section{Distribution and retrieval over OCI}
For artifact-only dispositions, the VEX document \rfckw{should} travel with the image it
describes. CycloneDX documents have a defined media type and can be stored as OCI
artifacts in the same registry as the image and associated with it through the registry's
referrers/attestation mechanism. A scanner that already pulls an image to detect
vulnerabilities can, in the same pass, resolve the image's attached VEX and apply the
dispositions automatically.

Deployment-context reachability attestations have a different subject: the service,
workload, or runtime component. They \rfckw{may} still be distributed as signed OCI
artifacts, but the reference is to a stable service/workload \code{bom-ref} and the
evidence binds the statement to the route and backend protocol configuration that
justified it. This allows the same surface attestation to be reused across scans and CVEs
while still forcing re-evaluation when the deployment surface changes. The result is the
same operational property as image-attached VEX: the \emph{a posteriori} judgment, once
made and signed, is rediscovered and reused without re-litigation, but without pretending
that deployment-context reachability is a property of the image alone.

\section{Optional interoperability: CycloneDX to CISA labels}
A FedRAMP-facing report needs nothing beyond the above. Some downstream consumers ---
agency tooling, or a CSAF advisory pipeline --- speak CISA's narrower
status-justification vocabulary instead. CISA, like CycloneDX, separates the two axes:
each CycloneDX \code{state} coarsens onto a CISA \emph{status} and each
\code{justification} onto a CISA \emph{justification label}, in both cases without loss
of safety (the mapping only ever generalizes; it never invents a stronger claim). This
crosswalk is \rfckw{optional} and exists solely for interchange.

\begin{table}[H]
\centering
\footnotesize
\renewcommand{\arraystretch}{1.25}
\begin{tabular}{>{\raggedright\arraybackslash}p{0.40\textwidth} >{\raggedright\arraybackslash}p{0.52\textwidth}}
\toprule
\textbf{CycloneDX justification / state} & \textbf{CISA label (coarsened)} \\
\midrule
\code{code\_not\_present} & \code{vulnerable\_code\_not\_present} \\
\code{code\_not\_reachable} & \code{vulnerable\_code\_not\_in\_execute\_path} \\
\code{requires\_configuration} & \code{vulnerable\_code\_not\_in\_execute\_path} \\
\code{requires\_dependency} & \code{vulnerable\_code\_not\_in\_execute\_path} \\
\code{requires\_environment} & \code{vulnerable\_code\_not\_in\_execute\_path} \\
\code{protected\_by\_compiler} & \code{inline\_mitigations\_already\_exist} \\
\code{protected\_at\_runtime} & \code{inline\_mitigations\_already\_exist} \\
\code{protected\_at\_perimeter} & \code{inline\_mitigations\_already\_exist} \\
\code{protected\_by\_mitigating\_control} & \code{inline\_mitigations\_already\_exist} \\
\midrule
\code{state: not\_affected} & \code{not\_affected} (paired with the justification above as the CISA label) \\
\code{state: false\_positive} & \code{not\_affected} (\code{component\_not\_present} where the component is truly absent) \\
\code{state: in\_triage} & \code{under\_investigation} \\
\code{state: exploitable} & \code{affected} \\
\code{state: resolved} / \code{resolved\_with\_pedigree} & \code{fixed} \\
\bottomrule
\end{tabular}
\caption{Optional coarsening of CycloneDX dispositions onto CISA's labels/statuses.}
\end{table}

\section{Standing and adoption}
CycloneDX introduced its vulnerability and VEX capabilities in version~1.4
(January~2022), ahead of CISA's status-justification guidance (June~2022); its richer
vocabulary is thus an independent, earlier design rather than a divergence from a later
one. CycloneDX is now an Ecma International standard (ECMA-424). Adoption is additive:
the disposition layer sits on top of the unchanged scoring pipeline of the method memo,
joined by CVE and component, and a provider \rfckw{may} emit the optional CISA crosswalk
for consumers that require it. As with the method's archetype catalog, this profile is
offered as a workable existence proof, not as a new FedRAMP standard --- the choice of
carrier remains the provider's.

\end{document}
